\clearpage
\chapter*{강연 182의 정오표}
\addcontentsline{toc}{chapter}{강연 182의 정오표}
\setcounter{footnote}{1}
\src{298}

\noindent\textbf{2. 형식기하학과 대수기하학.}

\begin{center}
[부르바키 세미나, 제11권, 1958/59, 제182호, 28쪽.]
\end{center}

\noindent\textbf{8쪽, 정리 6.} --- ``$f$가 준사영이다''라는 가정은
더 약한 ``$f$가 분리 사상이다''라는 가정으로 바꿀 수 있다. 이는 다음
결과에 따른 것이다(cf. SGA VIII, 6.2): 모든 준유한 분리 사상
$f:X\longrightarrow Y$는 준사영이다.

\src{299}
\noindent\textbf{12쪽 및 14쪽, 주 1.} --- J.-P.~세르는 표수가 $p>0$인
대수적으로 닫힌 체 $k$ 위의 차원 3인 비특이 사영 다양체를 구성했다. 이
다양체는 잉여체가 $k$이고 분수체의 표수가 0인 정역 국소환 위의 고유 스킴을
환원하여 얻어지는 것이 아니다.\footnote{Serre (Jean-Pierre). --- Exemples de
variétés projectives en caractéristique $p$ non relevables en caractéristique
zéro, \emph{Proc. Nat. Acad. Sc. U.S.A.}, t. 47, 1961, p. 108--109.} 멈퍼드도
비특이 사영 곡면인 이와 유사한 예를 찾았다고 한다.

\noindent\textbf{15쪽, 주 2와 3.} --- 나는 현재 여기서 추측한 결과들에
관해서는 덜 낙관적이다. 그러나 주 3의 말미에서 제기한 사영공간의 구조
변형에 관한 문제는 히로나카가 긍정적으로 해결했으며, 아벨 다양체에 관한
유사한 문제는 고이즈미가 해결했다.

\noindent\textbf{16쪽, 3행.} --- 멈퍼드는 최근 종수 $g$의 곡선들에 대한
모듈라이 스킴을 구성했다(cf. SHMI). 또한 정리 10은 모듈라이 이론의
``준위 $n$의 야코비 스킴''이 비특이임을 증명한다(더 나아가 이들은
$\mathbf Z$ 위에서 매끄럽기까지 하다).

\noindent\textbf{23쪽, 12행, 따름정리 5의 제1행.} --- 다음을 덧붙인다:
``$X$ 또는 $Y$가 $k$ 위에서 고유인 경우''.

§§ 1--5의 내용은 EGA III의 출판된 부분에 들어 있고, §§ 6과 7의 내용은
SGA III에 들어 있다. 기본군의 연구에 관해서는 SGA V, IX, X, XI 및
SGA 1962(강연 X, XII, XIII)를 보라. 여기에는 르프셰츠 유형의 정리들과
수많은 열린 문제가 다루어져 있다. 순하게 분기된 피복의 이론(cf. 정리 14)만
아직 정식으로 집필되지 않았다. 정리 14의 따름정리는 대수적으로 닫힌 체
위의 대수곡선에 대해, 그 표수와 서로소인 위수를 갖는 갈루아 피복을 완전히
결정하며, 다음 세 경우에 본질적으로 사용되었다.

1\textsuperscript{o} 임의의 표수에서 비특이 사영 곡면에 대한 피카르
부등식을 이구사가 증명할 때;

2\textsuperscript{o} 임의의 표수에서 일변수 함수체 위에 정의된 아벨
다양체를 구조군으로 하는 동차 주다발들의 군을 연구할 때(오그와
샤파레비치가 서로 독립적으로 발전시켰다);

3\textsuperscript{o} 대수다양체의 ``베유 코호몰로지''에 관한 몇몇 핵심
정리를 아르틴이 최근 증명할 때.
